\section{Обоснование цели и задач исследования}

На основе проведённого анализа существующих решений и выявленных пробелов в технологиях \texttt{Text-to-SQL} — в частности, отсутствия надёжных механизмов самокоррекции и ограничений при работе со сложными схемами данных — было сформулировано техническое обоснование цели и декомпозированных задач исследования.

\subsection{Обоснование цели работы}

Целью данной работы является разработка программного решения, которое устранит барьер между пользователями без технического образования и сложными реляционными базами данных. Создание полноценного \texttt{NLIDB}-интерфейса продиктовано необходимостью демократизации данных в крупных информационных системах, например, в сфере авиаперевозок. Это позволит конечным пользователям — менеджерам и аналитикам — получать оперативную информацию без ожидания выгрузок от ИТ-специалистов и тем самым снизит общую нагрузку на технический отдел предприятия.

Для достижения поставленной цели предстоит решить несколько взаимосвязанных задач.

Прежде всего, нужно проверсти оценку современных методов трансляции естественного языка в язык запросов. На сегодняшний день существует множество подходов и предобученных языковых моделей — как проприетарных, так и с открытым исходным кодом. Чтобы выбрать оптимальное «ядро» системы, потребуется провести сравнительный анализ их фичей.

Следующая задача — спроектировать \texttt{RAG}-архитектуру для работы с таблицами в схеме базы данных, с выявлением метаинформации и связей между таблицами. Передавать полную схему базы данных в контекстное окно языковой модели неэффективно и зачастую невозможно из-за ограничений на размер контекста. Архитектура \texttt{RAG} (Retrieval-Augmented Generation) позволяет динамически извлекать из векторного хранилища только те описания таблиц, полей и связей, которые действительно релевантны текущему вопросу пользователя. Такой подход гарантирует масштабируемость системы и сводит к минимуму риск галлюцинаций модели на сложных многотабличных базах данных.

Третья задача — разработка механизма самокоррекции ИИ-агента для повышения качества итогового ответа. Даже самые продвинутые языковые модели иногда генерируют \texttt{SQL}-запросы с синтаксическими или логическими ошибками. Подражая поведению человека-аналитика, ИИ-агент должен уметь перехватывать ошибки, возвращаемые СУБД, анализировать лог ошибки и автоматически переформулировать запрос. Внедрение такого цикла обратной связи (Self-Heal) заметно повысит автономность системы и долю успешно выполненных пользовательских запросов.

Четвёртая задача — провести оценку качества генерации после добавления в проект механизма самокоррекции. Любая дополнительная архитектурная надстройка требует валидации. Необходимо измерить метрики качества генерации, например точность выполнения запросов (Execution Accuracy), до и после внедрения механизма самокоррекции. Это позволит экспериментально подтвердить эффективность разработанного агентского подхода и обосновать его практическую ценность в рамках дипломного проекта.

Пятой задачей является сравнение существующих LLM в качестве \enquote{мозга системы}. Требуется протестировать отечественную разработку \texttt{Giga-chat 2} совместно с азиатскими моделями \texttt{Qwen3-Coder-Next} и \texttt{GLM-4.7}

Наконец, последней задачей является реализация микросервисной архитектуры. Интеллектуальный модуль должен обладать высокой отказоустойчивостью, гибкостью и независимостью компонентов. Разделение системы на микросервисы — сервис работы с \texttt{LLM}, сервис генерации эмбеддингов и векторного поиска, сервис-коннектор к реляционной базе данных — позволит независимо масштабировать наиболее нагруженные части системы, а также упростит интеграцию разработанного модуля в существующие корпоративные контуры предприятий.